% -------------------------------------------------------------------------
% WBS ID: RES-VER-BMK
% Deliverable: Analytical and Canonical Benchmark Suite
% Status: In progress
% Evidence status: Regional and local benchmark classes established
% Review status: Not reviewed
% Last updated: 2026-07-18
% Dependencies: RES-VER-FRM, RES-NUM-SPEC
% Downstream interfaces: RES-VER-MET, Software benchmark execution
% -------------------------------------------------------------------------

\section{RES-VER-BMK --- Analytical and Canonical Benchmark Suite}
\label{sec:res-ver-bmk}

\subsection{Purpose and Scope}
\label{sec:res-ver-bmk-purpose}

% TODO: Define what this Deliverable covers and explicitly excludes.

\subsection{Research Questions}
\label{sec:res-ver-bmk-research-questions}

% TODO: State the questions that must be answered before the Deliverable
% can be accepted.

\subsection{Terminology and Definitions}
\label{sec:res-ver-bmk-definitions}

% TODO: Define the technical terminology, symbols, quantities and
% classifications used in this Deliverable.

\subsection{Literature Evidence}
\label{sec:res-ver-bmk-literature-evidence}

% TODO: Record evidence from peer-reviewed papers, authoritative datasets,
% standards, technical reports and primary documentation.
%
% Do not create a detached literature-summary list. Explain how each source
% supports, contradicts or limits a project decision.

\subsubsection{Grilli et al. (2013)}
\label{sec:res-ver-bmk-evidence-grilli-2013}

\textcite{grilli2013tohoku} constructed a complete Tōhoku hindcast
linking a transient earthquake source, non-hydrostatic tsunami
generation, dispersive regional propagation, and nested coastal
inundation. The model was compared with near- and far-field
deep-ocean records, offshore GPS-wave gauges, run-up measurements, and
inundation observations.

The paper is retained as evidence for a stage-specific benchmark in
which source generation, offshore propagation, coastal transformation,
and inundation are assessed separately.

\subsubsection{Wei et al. (2013)}
\label{sec:res-ver-bmk-evidence-wei-2013}

\textcite{wei2013japan} used a source constrained by two deep-ocean
tsunami records to drive eleven nested inundation models along the
Japanese coastline. The resulting simulations were compared with
independent offshore GPS records, nearshore gauges, surveyed heights,
and mapped inundation areas.

The paper is retained as evidence that one regional source and
propagation solution can supply several local models, and that
observations used to infer the source must be distinguished from those
used to assess its predictive performance.

\subsection{Governing Relations and Quantitative Definitions}
\label{sec:res-ver-bmk-governing-relations}

% TODO: Record relevant equations, dimensionless groups, empirical models,
% extraction rules or quantitative definitions.
%
% Use align environments for displayed equations.
% Every displayed equation must have a unique \label{eq:...}.
% Do not place punctuation after displayed equations.

\subsection{Test Objective}
\label{sec:res-ver-bmk-test-objective}

% TODO: Complete this RES-VER Domain-specific subsection.

The benchmark suite shall verify solver capability before event
validation and loading validation are attempted. Local 3D benchmarks
shall isolate free-surface transport, pressure correction, wave
generation, wall loading and overtopping before the full
event-conditioned barrier model is used.

\subsection{Reference or Benchmark Solution}
\label{sec:res-ver-bmk-reference-benchmark-solution}

% TODO: Complete this RES-VER Domain-specific subsection.

The Tōhoku benchmark shall be divided into the following validation
stages:

\begin{enumerate}
  \item earthquake-source and seabed-deformation assessment;
  \item initial tsunami-generation assessment;
  \item near-source offshore propagation;
  \item regional propagation towards the selected coastline;
  \item nearshore transformation;
  \item run-up and inundation;
  \item local barrier interaction.
\end{enumerate}

Failure at one stage shall not be concealed by acceptable agreement at
a later stage. For example, an incorrect source shall not be accepted
solely on the basis of matching one inundation contour after roughness
or bathymetry has been adjusted.

The local V1 solver-capability benchmark classes are:

\begin{enumerate}
  \item still water with gravity and a flat free surface;
  \item VOF advection of a prescribed interface;
  \item dam-break or bore propagation;
  \item regular or solitary wave generation and absorption in a
    numerical wave tank;
  \item solitary-wave run-up on a slope;
  \item wave impact on a vertical wall or simple fixed obstacle;
  \item overtopping of a simple impermeable barrier;
  \item flow around a fixed object with pressure and force extraction;
  \item mesh and timestep refinement for the local quantities of
    interest.
\end{enumerate}

\textcite{watanabe2022validation} and \textcite{qin2018comparison}
support the separation between hydrodynamic validation and impact/load
validation. Source role: `3D-VALIDATION`. Project conclusion:
pressure, velocity and force benchmarks are required before local
barrier-performance claims are accepted.

\subsection{Test Configuration}
\label{sec:res-ver-bmk-test-configuration}

% TODO: Complete this RES-VER Domain-specific subsection.

\subsection{Simulation Duration}
\label{sec:res-ver-bmk-simulation}

The benchmark interval shall include all waves that materially affect
the selected coastal and structural quantities.

The local forcing history shall not be reduced automatically to:

\begin{itemize}
  \item the first arriving crest;
  \item the maximum isolated amplitude;
  \item one idealised solitary wave;
  \item a short interval surrounding initial impact.
\end{itemize}

Later waves may propagate over residual inundation water, interact with
coastal reflections, and produce the maximum inundation, overtopping,
or cumulative load after the first arrival.

\subsection{Measured Quantities}
\label{sec:res-ver-bmk-measured-quantities}

% TODO: Complete this RES-VER Domain-specific subsection.

Local benchmark outputs shall include phase-volume conservation,
free-surface elevation, wave-front arrival, depth-dependent velocity,
run-up, overtopping volume, pressure histories, pressure distribution,
resultant force, impulse, moment and reflected/transmitted wave
measures.

\subsection{Error Metrics}
\label{sec:res-ver-bmk-error-metrics}

% TODO: Complete this RES-VER Domain-specific subsection.

\subsection{Tolerance and Acceptance Criteria}
\label{sec:res-ver-bmk-tolerance-acceptance-criteria}

% TODO: Complete this RES-VER Domain-specific subsection.

\subsection{Execution Handoff}
\label{sec:res-ver-bmk-execution-handoff}

% TODO: Complete this RES-VER Domain-specific subsection.

\subsection{Candidate Approaches}
\label{sec:res-ver-bmk-candidate-approaches}

% TODO: Define the credible candidate approaches considered.

\subsection{Critical Comparison}
\label{sec:res-ver-bmk-critical-comparison}

% TODO: Compare candidates using physically and project-relevant criteria.
% Do not select an approach solely on popularity or implementation ease.

\subsection{Assumptions and Applicability}
\label{sec:res-ver-bmk-assumptions}

% TODO: State assumptions and the conditions under which they are valid.

\subsection{Limitations and Failure Conditions}
\label{sec:res-ver-bmk-limitations}

% TODO: State where the evidence, model or selected approach ceases to be
% reliable or applicable.

\subsection{Project-Specific Conclusions}
\label{sec:res-ver-bmk-conclusions}

% TODO: Convert the evidence into conclusions specific to the Studentship
% project. Do not merely restate literature findings.

\subsection{Selected Approach and Rationale}
\label{sec:res-ver-bmk-selected-approach}

% TODO: Record the selected approach only after sufficient evidence exists.
% State the alternatives rejected and the reasons for rejection.

The selected benchmark approach is progressive. A full
event-conditioned local barrier simulation shall not be the first test
of the 3D solver; simpler canonical cases shall verify each numerical
and physical mechanism before combined validation.

\subsection{Unresolved Decisions}
\label{sec:res-ver-bmk-unresolved-decisions}

% TODO: Record unresolved questions, required evidence and decision owners.

\subsection{Requirements and Handoffs}
\label{sec:res-ver-bmk-requirements}

% TODO: State requirements passed to other Research Domains, Software,
% verification activities or later execution work.

Software benchmark reports shall identify the benchmark class,
geometry, mesh, timestep limits, boundary conditions, turbulence mode,
measured quantities and pass/fail criteria. `RES-VER-MET` shall define
the numerical values of the acceptance metrics.

\subsection{Deliverable Completion Record}
\label{sec:res-ver-bmk-completion-record}

% TODO: Record whether the Deliverable is:
%
% - Not started
% - In progress
% - Draft complete
% - Reviewed
% - Accepted
%
% Acceptance requires:
%
% - the research questions to be answered;
% - claims to be supported by appropriate evidence;
% - assumptions and limitations to be explicit;
% - the project-specific conclusion to be stated;
% - downstream requirements to be traceable;
% - unresolved decisions to be closed or formally deferred.
